\documentclass[11pt]{article}
\usepackage[T1]{fontenc}
\usepackage{lmodern,microtype,amsmath,amssymb,amsthm,mathtools,mathrsfs,tikz-cd}
\usepackage[margin=1in]{geometry}
\usepackage{xurl,hyperref}
\hypersetup{colorlinks=true,linkcolor=blue,citecolor=blue,urlcolor=blue,
  pdftitle={Deduction and essential incompleteness in recursion categories},
  pdfauthor={Florian Lengyel}}
\newtheorem{theorem}{Theorem}[section]
\newtheorem{proposition}[theorem]{Proposition}
\newtheorem{lemma}[theorem]{Lemma}
\newtheorem{corollary}[theorem]{Corollary}
\theoremstyle{definition}
\newtheorem{definition}[theorem]{Definition}
\newtheorem{example}[theorem]{Example}
\newcommand{\R}{\mathscr R}
\newcommand{\Pos}{\mathbf{Pos}}
\newcommand{\BDLat}{\mathbf{BDLat}}
\newcommand{\Tot}{\operatorname{Tot}}
\newcommand{\ran}{\operatorname{ran}}
\newcommand{\Th}{\mathsf{Th}}
\newcommand{\Cons}{\mathsf{Cons}}
\newcommand{\id}{\operatorname{id}}
\title{Deduction and essential incompleteness\protect\\in recursion categories}
\author{Florian Lengyel}
\date{29 October 2026}
\begin{document}
\maketitle

\begin{abstract}
Inference rules in a recursion category act on restriction idempotents by
taking ranges. Stable unions of the increasing deduction stages yield the
least domain closed under the rules, and hence a closure operator whose fixed
points are theories. A representation condition on the positive and negative
domains of a theory implies their computational inseparability. Consequently,
every consistent extension represented by a closed domain is incomplete.
In the category of partial recursive maps, the required unions exist because
the deduction stages are uniformly computably enumerable. In the arithmetic
instance, every consistent computably enumerable extension of a consistent
theory strongly defining all recursive sets numeralwise is incomplete; the
base theory need not be computably enumerable.
\end{abstract}

\section{Domains and ranges}

Deduction requires both operations producing conclusions and a means of
collecting their outputs. We express the operations as partial morphisms and
the collections as restriction idempotents. Ranges provide direct images of
these domains, and Heller's stable unions provide closure under finite
iterations of the rules \cite{Heller}.

We write $gf=g\circ f$. Let $\R$ be a locally small $P$-category with zero
morphisms, in the sense of \cite{Heller}. We use the restriction notation
of \cite{RestrictionI},
written $f\mapsto\bar f$, and its binary restriction products, written
$X\boxtimes Y$. The restriction of $f:X\to Y$ is an endomorphism
$\bar f:X\to X$. The restriction identities are
\[
 f\bar f=f,\qquad
 \bar f\bar g=\bar g\bar f,\qquad
 \overline{g\bar f}=\bar g\bar f,\qquad
 \bar h f=f\overline{hf}.
\]
In the middle two identities $f$ and $g$ have a common source; in the last,
$f:X\to Y$ and $h:Y\to Z$. We assume $\overline{0_{X,Y}}=0_{X,X}$.
A morphism is \emph{total} when $\bar f=1_X$. The projections of a restriction
product are total, and the pairing of $f:Z\to X$ and $g:Z\to Y$ satisfies
\[
 \pi_1\langle f,g\rangle=f\bar g,\qquad
 \pi_2\langle f,g\rangle=g\bar f,\qquad
 \overline{\langle f,g\rangle}=\bar f\bar g.
\]
Pairing is determined by these equations. The near product on morphisms
preserves composition and satisfies
$\overline{f\boxtimes g}=\bar f\boxtimes\bar g$.
The total morphisms form a category $\Tot(\R)$ with binary products.
Its diagonal is $\Delta_X=\langle1_X,1_X\rangle$.

A \emph{domain} on $X$ is a restriction idempotent $e:X\to X$, meaning
$\bar e=e$. Write
\[
 \operatorname{Dom}(X)=\{e:X\to X\mid\bar e=e\}.
\]
The order $e\le d$ is defined by $de=e$. Domains on the same object commute,
their meet is $ed$, and their least and greatest elements are $0_X$ and $1_X$.
For a fixed domain $q$ on $U$, write
$E_q=\{a\in\operatorname{Dom}(U):a\le q\}$.

The intended examples are recursion categories: prodominical $P$-categories
whose objects are all isomorphic and which admit a Turing morphism
\cite{DPH,Heller}. On a chosen object $U$, the latter is a morphism
$\tau:U\boxtimes U\to U$ such that, for every $f:X\boxtimes U\to U$,
there is a total $g:X\to U$ with $f=\tau(g\boxtimes1_U)$.
The constructions below are stated with the particular domain, range, and
union hypotheses they require.

\begin{definition}[Range]\label{def:range}
A \emph{range} of $f:X\to Y$ is a domain $r\in\operatorname{Dom}(Y)$ such that $rf=f$ and,
for every object $Z$ and every pair $\beta,\beta':Y\to Z$,
\[
 \beta f=\beta'f\quad\Longrightarrow\quad\beta r=\beta'r.
\]
The category \emph{has ranges} if every morphism has a range.
\end{definition}

This is the range definition of Di Paola--Heller \cite[Definition 5.1]{DPH}
and Heller \cite{Heller}. A range is unique. Indeed, if $r,r'$ are ranges of
$f$, their universal properties give $r'r=r$ and $rr'=r'$, and domains
commute. We write the unique range as $\ran(f)$.

\begin{lemma}\label{lem:range}
In a $P$-category with ranges, for $f:X\to Y$, $g:Y\to Z$, and domains
$e\in\operatorname{Dom}(X)$ and $d\in\operatorname{Dom}(Y)$,
\begin{align}
 \ran(f)\le d&\quad\Longleftrightarrow\quad df=f,\label{eq:least-range}\\
 \ran(e)&=e,\label{eq:range-domain}\\
 \ran(gf)&=\ran(g\ran(f)),\label{eq:range-composition}\\
 \ran(df)&=d\ran(f).\label{eq:range-meet}
\end{align}
\end{lemma}
\begin{proof}
If $df=f$, apply the universal property to $d$ and $1_Y$ to obtain
$d\ran(f)=\ran(f)$. The converse follows by composing this equality with
$f$. A domain $e$ satisfies the range conditions for itself, proving
\eqref{eq:range-domain}.

To prove \eqref{eq:range-composition}, the domain $r=\ran(g\ran(f))$
receives $gf$. If $\beta gf=\beta'gf$, the range property for $f$ first gives
$\beta g\ran(f)=\beta'g\ran(f)$; the property for $g\ran(f)$ then gives
$\beta r=\beta'r$. Thus $r$ is a range of $gf$.

For \eqref{eq:range-meet}, the domain $d\ran(f)$ receives $df$, since $d$
commutes with $\ran(f)$. Agreement after $df$, followed by the range
property for $f$ applied to $\beta d$ and $\beta'd$, gives agreement after
$d\ran(f)$. Uniqueness proves the identity.
\end{proof}

Ranges define a covariant functor
\[
 E_!:\R\longrightarrow\Pos.
\]
On objects it assigns $E_!(X)=\operatorname{Dom}(X)$. On $f:X\to Y$ it assigns
\[
 f_!:\operatorname{Dom}(X)\longrightarrow\operatorname{Dom}(Y),\qquad f_!(e)=\ran(fe).
\]
If $e\le e'$, then $\ran(fe')$ receives $fe$, proving monotonicity.
Equations \eqref{eq:range-domain} and \eqref{eq:range-composition} give
$ (1_X)_!=1_{\operatorname{Dom}(X)}$ and $(gf)_!=g_!f_!$.

For any $f:X\to Y$, define the inverse-image map
\[
 f^*:\operatorname{Dom}(Y)\longrightarrow\operatorname{Dom}(X),\qquad f^*(d)=\overline{df}.
\]
The restriction identities give preservation of meets and
$(gf)^*=f^*g^*$. They also give the precise relation
\begin{equation}\label{eq:local-adjunction}
 f_!(e)\le d\quad\Longleftrightarrow\quad
 e\bar f\le f^*(d).
\end{equation}
Indeed, the left side is equivalent to $dfe=fe$. Restricting this equality
gives $\overline{df}e=\bar f e$; conversely that equality and
$df=f\overline{df}$ give $dfe=fe$.
For total $f$ this is the adjunction $f_!\dashv f^*$.

\section{Stable unions and inverse image}

We use Heller's notion of stable union \cite{Heller}. A finite family can be
regarded as a countable family by repetition.

\begin{definition}[Stable union]\label{def:stable-union}
Let $(e_i)_{i\in I}$ be a nonempty finite or countable family in $\operatorname{Dom}(X)$.
A domain $e\in\operatorname{Dom}(X)$ is its \emph{union} if $e_i\le e$ for every $i$ and,
for every parallel pair $\beta,\beta':X\to Y$,
\[
 (\forall i\in I,\ \beta e_i=\beta'e_i)
 \quad\Longrightarrow\quad\beta e=\beta'e.
\]
The union is \emph{stable} if, for every object $W$, the domain
$1_W\boxtimes e$ is the union of the $1_W\boxtimes e_i$.
\end{definition}

Every such union is a least upper bound: apply its universal property to
an upper bound $d$ and $1_X$. Thus it is unique, and we write it as
$\bigvee_i e_i$. Symmetry of the near product gives stability in the other
coordinate as well. The empty union is $0_X$.

\begin{lemma}\label{lem:union-calculus}
Let $\R$ be a $P$-category with ranges.
\begin{enumerate}
\item If $e=\bigvee_i e_i$ is a union and $d\in\operatorname{Dom}(X)$, then $de$ is
      the union of the $de_i$.
\item Direct image preserves unions:
      $f_!(\bigvee_i e_i)=\bigvee_i f_!(e_i)$.
\item Inverse image along a total morphism preserves stable unions.
\end{enumerate}
\end{lemma}
\begin{proof}
For the first assertion, suppose $\beta de_i=\beta'de_i$ for every $i$.
Apply the union property to $\beta d$ and $\beta'd$. The result is agreement
after $de$; the upper-bound condition follows from monotonicity of meet.

For the second, suppose $\beta\ran(fe_i)=\beta'\ran(fe_i)$ for every $i$.
Composing with $fe_i$ gives $\beta fe_i=\beta'fe_i$. The union property for
$e$, and then the range property for $fe$, give agreement after $\ran(fe)$.
Monotonicity of $f_!$ gives the upper-bound condition.

For the third, let $h:X\to Y$ be total, and let $d=\bigvee_i d_i$ be a
stable union in $\operatorname{Dom}(Y)$. Put $g=\langle1_X,h\rangle:X\to X\boxtimes Y$ and
$R=\ran(g)$. The restriction-product identities and
\eqref{eq:range-meet} give, for each domain $v\in\operatorname{Dom}(Y)$,
\[
 (1_X\boxtimes v)g=g\,h^*v,
 \qquad
 (1_X\boxtimes v)R=\ran(g\,h^*v).
\]
Stability and the first assertion show that $(1_X\boxtimes d)R$ is the union
of $(1_X\boxtimes d_i)R$. Apply the second assertion to the projection
$\pi_1:X\boxtimes Y\to X$. Since $\pi_1g=1_X$,
\[
 (\pi_1)_!\bigl((1_X\boxtimes v)R\bigr)
 =\ran(\pi_1g\,h^*v)=h^*v.
\]
This proves that $h^*d$ is the union of the $h^*d_i$.
For each object $W$, apply the same argument to $1_W\boxtimes h$ and
the stable union of the $1_W\boxtimes d_i$. Iterated near products preserve
this union by stability and associativity. Since
$(1_W\boxtimes h)^*(1_W\boxtimes v)=1_W\boxtimes h^*v$, the union of
the inverse images is itself stable.
\end{proof}

Assume henceforth that finite stable unions exist. Lemma~\ref{lem:union-calculus}
shows that each $\operatorname{Dom}(X)$ is a bounded distributive lattice. Let $\BDLat$ be
the category of such lattices and homomorphisms preserving both bounds,
finite meets, and finite joins. The contravariant domain functor is
\[
 E:\Tot(\R)^{\mathrm{op}}\longrightarrow\BDLat.
\]
On objects it assigns $E(X)=\operatorname{Dom}(X)$. On $f^{\mathrm{op}}:Y\to X$, corresponding
to a total $f:X\to Y$, it assigns $f^*:E(Y)\to E(X)$. The restriction
identities give the identity and composition laws and preservation of meets;
totality and restriction zeros give preservation of the bounds. Preservation
of joins is Lemma~\ref{lem:union-calculus}.

\begin{lemma}\label{lem:chain-powers}
If $e_0\le e_1\le\cdots$ has stable union $e$, then, for every $n\ge1$,
$e^{\boxtimes n}$ is the union of $e_k^{\boxtimes n}$. For every object $P$,
$1_P\boxtimes e^{\boxtimes n}$ is the union of
$1_P\boxtimes e_k^{\boxtimes n}$.
\end{lemma}
\begin{proof}
Suppose two parallel morphisms agree after every $e_k^{\boxtimes n}$.
They agree after each mixed product
$e_{k_1}\boxtimes\cdots\boxtimes e_{k_n}$, since that product is below
$e_k^{\boxtimes n}$ for $k=\max(k_1,\ldots,k_n)$. Apply stability successively
in each coordinate. At each application, compose the two morphisms with the
fixed domains in the remaining coordinates. This gives agreement after
$e^{\boxtimes n}$. All members lie below that domain. The same argument
with the additional fixed factor $1_P$ proves the second assertion.
\end{proof}

\section{Deductive closure and theories}

Fix an object $U$, a domain $q\in E(U)$, and a domain $\ell\le q$ of fixed
axioms. A rule with $n_i\ge1$ premises and parameter object $P_i$ is a
morphism
\[
 r_i:P_i\boxtimes U^{\boxtimes n_i}\longrightarrow U
\]
satisfying
\begin{equation}\label{eq:rule-support}
 r_i=q\,r_i(1_{P_i}\boxtimes q^{\boxtimes n_i}).
\end{equation}
A rule without parameters has type $U^{\boxtimes n_i}\to U$, with the
corresponding support equation. Rules with no premises are included in
$\ell$. Fix finitely many rules. Each defines a monotone map
\[
 R_i:E_q\longrightarrow E_q,\qquad
 R_i(a)=\ran\bigl(r_i(1_{P_i}\boxtimes a^{\boxtimes n_i})\bigr),
\]
with the parameter factor omitted for a rule without parameters.
Equation~\eqref{eq:rule-support} puts the image below $q$.

A domain $t\le q$ is \emph{closed under the rules} if $R_i(t)\le t$ for
each $i$. By \eqref{eq:least-range}, this condition is equivalent to
\[
 t\,r_i(1_{P_i}\boxtimes t^{\boxtimes n_i})
 =r_i(1_{P_i}\boxtimes t^{\boxtimes n_i}).
\]
For $a\in E_q$, its deduction stages are
\begin{equation}\label{eq:stages}
 a_0=\ell\vee a,\qquad
 a_{k+1}=a_k\vee\bigvee_i R_i(a_k).
\end{equation}
These form an increasing sequence in $E_q$.

\begin{theorem}[Deductive closure]\label{thm:closure}
Let $\R$ be a $P$-category with ranges and finite stable unions. Fix the
domain $q$, fixed axioms $\ell\le q$, and finitely many rules satisfying
\eqref{eq:rule-support}. Assume that, for every $a\le q$, the sequence
\eqref{eq:stages} has a stable union. There is a closure functor
\[
 c_q:E_q\longrightarrow E_q.
\]
On objects it is defined by $c_q(a)=\bigvee_k a_k$; an inequality $a\le b$
is sent to $c_q(a)\le c_q(b)$. Its value is the least rule-closed domain
containing $\ell\vee a$. In particular, $a\le c_q(a)$ and $c_q^2=c_q$.
\end{theorem}
\begin{proof}
All stages lie below $q$, so their union does as well. Lemmas
\ref{lem:union-calculus} and \ref{lem:chain-powers} imply
\[
 R_i(c_q(a))=\bigvee_k R_i(a_k)
 \le\bigvee_k a_{k+1}=c_q(a).
\]
Thus $c_q(a)$ is rule-closed. If $t$ is rule-closed and $\ell\vee a\le t$,
induction gives $a_k\le t$ for every $k$, hence $c_q(a)\le t$.

If $a\le b$, monotonicity of the rules gives the corresponding inequality
at every stage, and therefore $c_q(a)\le c_q(b)$. These inequalities define
the functor on the poset category. Extensivity follows from the initial
stage. Leastness and closedness give $c_q(c_q(a))=c_q(a)$.
\end{proof}

Some applications distinguish a subdomain in which conclusions are selected.
Fix $s\le q$. Define the functor
\[
 c_s:E_s\longrightarrow E_s
\]
on domains by $c_s(a)=s\,c_q(a)$ and on inequalities by monotonicity.
It is extensive because $a\le s$ and $a\le c_q(a)$. Moreover,
\[
 c_s(c_s(a))=s\,c_q(s\,c_q(a))
 \le s\,c_q(c_q(a))=s\,c_q(a)=c_s(a).
\]
Extensivity gives the reverse inequality, so $c_s$ is a closure operator.

\begin{definition}[Theory and extension]\label{def:theory}
The category $\Th_c$ is the full subposet of $E_s$ consisting of domains
$t$ with $c_s(t)=t$. Its objects are called \emph{theories}. An extension
of $b\in\Th_c$ is a theory $t$ with $b\le t$.
\end{definition}

The inclusion functor $i:\Th_c\hookrightarrow E_s$ sends each domain and
inequality to itself. Define the functor
\[
 L:E_s\longrightarrow\Th_c
\]
on objects by $L(a)=c_s(a)$ and on morphisms by $a\le b\mapsto c_s(a)\le c_s(b)$.
It is left adjoint to $i$, since for closed $t$,
\begin{equation}\label{eq:reflection}
 L(a)\le t\quad\Longleftrightarrow\quad a\le i(t).
\end{equation}

\subsection{Morphisms of presentations}

Fix the ambient category $\R$. A presentation consists of the data
$(U,q,s,\ell,(r_i))$ of Theorem~\ref{thm:closure}, including the specified
stable unions of its deduction chains. Define a category $\Cons$ as follows.
Its objects are these presentations. A morphism from a presentation with
selected domain $s$ to one with selected domain $s'$ is a morphism
$f:U\to U'$ of $\R$ satisfying
\[
 \bar f=s,\qquad s'f=f,
 \qquad c_s f^*(a)\le f^*c_{s'}(a)\quad(a\le s').
\]
The inverse image here has type $f^*:E_{s'}\to E_s$. The identity is $s$;
composition is composition in $\R$. The restriction identities give the
support equations for a composite, and the displayed inequalities compose
using monotonicity and $(gf)^*=f^*g^*$.

There is a contravariant functor
\[
 \Th:\Cons^{\mathrm{op}}\longrightarrow\Pos.
\]
It assigns to a presentation its poset $\Th_c$, and to $f^{\mathrm{op}}$
the restricted inverse-image map $f^*$. If $t$ is closed, continuity gives
$c_s(f^*t)\le f^*c_{s'}(t)=f^*t$, and extensivity gives equality. Thus the
map lands in the closed domains; its functor laws follow from inverse image.
The analogous functor $\mathcal E:\Cons^{\mathrm{op}}\to\Pos$ assigns
$E_s$ to a presentation and $f^*$ to a morphism. The inclusions form a natural
transformation $i:\Th\Rightarrow\mathcal E$, because both routes in the
naturality square are the same restricted inverse-image map. The closure
maps satisfy the stated inequality; equality is an additional condition.

\section{Representation and essential incompleteness}

Fix a presentation as above, with selected domain $s\in E(U)$ and closure
$c_s$. Supply a morphism
\begin{equation}\label{eq:negative-operation}
 \nu:U\to U,\qquad\bar\nu=s,\qquad s\nu=\nu.
\end{equation}
For $t\le s$, define its \emph{negative domain} by
$t^-=\nu^*t=\overline{t\nu}$. This is a domain below $s$.
A theory $t$ is \emph{consistent} if $t\wedge t^-=0_U$, and
\emph{complete} if $t\vee t^-=s$. It is \emph{incomplete} if the latter
equality fails.

A pair $(e_+,e_-)$ in $E(U)$ is \emph{complementary} if its meet is $0_U$
and its join is $1_U$. A pair $(d_+,d_-)$ in $E_s$ is complementary
\emph{relative to $s$} if its meet is $0_U$ and its join is $s$.
Such a pair \emph{separates} $p_+,p_-\le s$ if $p_+\le d_+$ and
$p_-\le d_-$. The domains $p_+,p_-$ are \emph{computationally inseparable}
if they are disjoint and have no such separator.

Fix a total morphism
\[
 \sigma:U\boxtimes U\longrightarrow U,\qquad s\sigma=\sigma.
\]
For each total $g:U\to U$, put $f_g=\sigma(g\boxtimes1_U)$.

\begin{definition}[Representation]\label{def:representation}
A consistent theory $b\in\Th_c$ satisfies the \emph{representation condition}
for $\sigma$ and $\nu$ if, for every complementary pair $(e_+,e_-)$ in
$E(U)$, there exists a total $g:U\to U$ such that
\begin{equation}\label{eq:representation}
 f_g^*b=\pi_2^*e_+,\qquad
 f_g^*b^-=\pi_2^*e_-.
\end{equation}
\end{definition}

\begin{theorem}[Computational inseparability]\label{thm:inseparability}
Let $\R$ be a $P$-category with ranges and finite stable unions, and let
$U$ satisfy $0_U\ne1_U$. Fix a domain $s$ on $U$, a closure $c_s:E_s\to E_s$,
and morphisms $\nu:U\to U$ and $\sigma:U\boxtimes U\to U$ satisfying
$\bar\nu=s$, $s\nu=\nu$, $\bar\sigma=1_{U\boxtimes U}$, and
$s\sigma=\sigma$. If a consistent theory $b$ satisfies the representation
condition, then $b$ and $b^-$ are computationally inseparable.
\end{theorem}
\begin{proof}
Suppose that $(d_+,d_-)$ separates $(b,b^-)$ and is complementary relative to
$s$. Put $\delta=\sigma\Delta_U:U\to U$. This is total and $s\delta=\delta$.
The pair
\[
 e_+=\delta^*d_-,\qquad e_-=\delta^*d_+
\]
is complementary in $E(U)$: inverse image preserves finite meets and joins,
and $\delta^*s=1_U$. Choose $g$ representing this pair. Set
$\gamma=\langle1_U,g\rangle$ and $h=\delta g$. The restriction-product
identities give $f_g\gamma=h$ and $\pi_2\gamma=g$, yielding the commutative
square
\[
\begin{tikzcd}[column sep=large]
 U \arrow[r,"\gamma"] \arrow[d,"g"'] & U\boxtimes U \arrow[d,"f_g"]\\
 U \arrow[r,"\delta"'] & U.
\end{tikzcd}
\]
Pulling \eqref{eq:representation} back along $\gamma$ gives
\[
 h^*b=h^*d_-,\qquad h^*b^-=h^*d_+.
\]
The separator inequalities therefore imply both
$h^*d_-\le h^*d_+$ and $h^*d_+\le h^*d_-$. These two domains are
complementary, since $h$ is total and $sh=h$. Their equality forces
$0_U=1_U$, a contradiction.
\end{proof}

\begin{theorem}[Essential incompleteness]\label{thm:essential}
Let $\R$ have ranges and finite stable unions, and let a presentation satisfy
the deduction-chain hypotheses of Theorem~\ref{thm:closure}. Suppose
$0_U\ne1_U$, and fix $\nu$ and $\sigma$ as in
Theorem~\ref{thm:inseparability}. If a consistent theory $b\in\Th_c$
satisfies the representation condition, every consistent theory
$t\in\Th_c$ with $b\le t$ is incomplete.
\end{theorem}
\begin{proof}
Monotonicity gives $b^-=\nu^*b\le\nu^*t=t^-$. If $t$ were complete,
consistency and completeness would make $(t,t^-)$ a complementary pair
relative to $s$ separating $(b,b^-)$. This contradicts
Theorem~\ref{thm:inseparability}.
\end{proof}

The representation condition is inherited by consistent extensions. To see
this, let $g$ represent $(e_+,e_-)$ in $b$, and suppose $b\le t$ with $t$
consistent. Write $x=f_g^*t$ and $y=f_g^*t^-$. Then
\[
 \pi_2^*e_+\le x,\qquad\pi_2^*e_-\le y,\qquad x\wedge y=0_{U\boxtimes U}.
\]
Since the two lower bounds are complementary, distributivity gives
$x=\pi_2^*e_+$ and $y=\pi_2^*e_-$. Thus the same $g$ represents the pair
in $t$.

\section{The arithmetic realization}

Let $\R_{\mathrm{rec}}$ have objects $\mathbb N^k$, where $k\ge1$, and partial
recursive maps as morphisms, with equality of partial functions. Composition
is ordinary composition of partial maps. The near product is the product on
the underlying sets and on partial maps. Recursive pairing and its inverses
identify all the objects up to isomorphism, and a universal partial recursive
function supplies a Turing morphism. This is a numerical recursion category
\cite{DPH,Heller}.

For a c.e. subset $A\subseteq\mathbb N^k$, let $1_A$ denote the partial
identity with domain $A$. The domains of $\R_{\mathrm{rec}}$ are exactly
these partial identities. Their order is inclusion of the corresponding
sets; finite meets and joins are intersection and union. For a partial
recursive map $f$, its image is c.e., and
\[
 \ran(f)=1_{\operatorname{im}(f)}.
\]
Indeed, agreement after $f$ implies agreement, including definedness, at
every point of its image. This proves the universal range condition.

\begin{proposition}\label{prop:effective-union}
Fix $r\ge1$. Suppose $(A_k)_{k\in\mathbb N}$ is a uniformly c.e. family of
subsets of $\mathbb N^r$, meaning that
$\{(k,x):x\in A_k\}$ is c.e. Then the domains $1_{A_k}$ have the stable
union $1_{\bigcup_k A_k}$ in $\R_{\mathrm{rec}}$.
\end{proposition}
\begin{proof}
Let $D=\{(k,x):x\in A_k\}\subseteq\mathbb N^{r+1}$ and let
$\delta=1_D:\mathbb N^{r+1}\rightharpoonup\mathbb N^{r+1}$. For the total
projection $\pi_2:\mathbb N^{r+1}\to\mathbb N^r$ onto the $x$-coordinates,
\[
 \ran(\pi_2\delta)=1_{\bigcup_k A_k}.
\]
Thus the proposed union is a morphism of $\R_{\mathrm{rec}}$.
If $\beta1_{A_k}=\beta'1_{A_k}$ for every $k$, then for every
$x\in\bigcup_k A_k$ a stage containing $x$ shows that $\beta(x)$ and
$\beta'(x)$ have the same definedness and value. Outside the union both
composites with its partial identity are undefined. This proves the union
condition. For every object $W$,
\[
 W\times\bigcup_k A_k=\bigcup_k(W\times A_k),
\]
and the same argument proves stability.
\end{proof}

Arbitrary increasing countable families need not have unions. If
$B\subseteq\mathbb N$ is not c.e., each finite set
$B_k=B\cap\{0,\ldots,k\}$ is recursive, but its increasing sequence has no
union in the domain poset. Any upper bound contains $B$, while for each
$x\notin B$ the recursive set $\mathbb N\setminus\{x\}$ is an upper bound.
A least upper bound would therefore equal $B$.

\subsection{Effective deduction}

Fix a computable coding of the terms and formulas of the language of
arithmetic, with decidable recognition of well-formed formulas and sentences.
Let $F$ and $S$ be their respective sets of codes, and put $q=1_F$ and
$s=1_S$. Choose a Hilbert calculus for classical first-order logic with
equality and a recursive set $\Lambda$ of logical axiom instances; put
$\ell=1_\Lambda$. Nonlogical axioms are sentences.

The calculus has the partial recursive rule maps
\[
 \mathrm{mp}:\mathbb N^2\rightharpoonup\mathbb N,
 \qquad \mathrm{gen}:\mathbb N^2\rightharpoonup\mathbb N,
\]
where
\[
 \mathrm{mp}(\ulcorner\varphi\urcorner,\ulcorner\varphi\to\psi\urcorner)
 =\ulcorner\psi\urcorner,
 \qquad
 \mathrm{gen}(v,\ulcorner\varphi\urcorner)=\ulcorner\forall v\,\varphi\urcorner.
\]
The first map is defined exactly on instances of modus ponens. The second
checks that its second input is a formula code; its first input indexes a
variable. Thus $\mathrm{mp}$ is a binary rule and $\mathrm{gen}$ a unary rule with parameter
object $\mathbb N$. Generalization is unrestricted because the nonlogical
axioms are sentences. Formula domains are used for intermediate deductions;
the selected domain $s$ restricts their conclusions to sentences.

For every c.e. domain $a\le q$, the stage construction becomes
\begin{align*}
 a_0&=\ell\vee a,\\
 a_{k+1}&=a_k\vee\ran\bigl(\mathrm{mp}(a_k\boxtimes a_k)\bigr)
                  \vee\ran\bigl(\mathrm{gen}(1_{\mathbb N}\boxtimes a_k)\bigr).
\end{align*}
These stages are uniformly c.e. To verify uniformity, fix an effective
numbering of the c.e. sets. From an index for $\operatorname{dom}(a_k)$,
one effectively obtains indices for its finite products, their images under
the two fixed rule maps, and the finite union displayed above. A total
recursive transformation of indices therefore takes an index for one stage
to an index for the next. Iterating this transformation computes a stage
index uniformly in $k$. A universal simulation consequently enumerates the
relation $(k,x)$ with $x\in\operatorname{dom}(a_k)$.
Proposition~\ref{prop:effective-union} supplies exactly the stable unions
required by Theorem~\ref{thm:closure}.

For a c.e. set $A\subseteq S$ of nonlogical axioms, induction on proof trees
shows that the union of the stages starting at $\ell\vee1_A$ is the domain
of all formula conclusions derivable from $A$. Each stage adds one rule
application to earlier conclusions. Conversely, every finite proof tree
has a finite height and hence its conclusion appears at a finite stage.
It follows that
\begin{equation}\label{eq:arithmetic-closure}
 c_s(1_A)=1_{\{\ulcorner\theta\urcorner:\theta\text{ is a sentence and }A\vdash\theta\}}.
\end{equation}
Thus the fixed points of $c_s$ are exactly the c.e. deductively closed sets
of sentences. This includes the theorem sets of Robinson's $Q$, Peano
arithmetic, and their c.e. axiomatic extensions.

\subsection{Representation and arbitrary base theories}

Define $\nu:\mathbb N\rightharpoonup\mathbb N$ by
$\nu(\ulcorner\theta\urcorner)=\ulcorner\neg\theta\urcorner$ on sentence
codes, leaving it undefined elsewhere. Then $\bar\nu=s$ and $s\nu=\nu$.
For a theorem domain $t$, the domain $t^-$ is precisely the set of sentence
codes refutable in that theory. The categorical consistency and completeness
conditions therefore have their usual syntactic meanings.

Choose a total computable enumeration $u\mapsto\varphi_u(v_0)$ of the
formulas whose only possible free variable is $v_0$, covering every such
formula. For example, use a syntactic code when it has this property and a
fixed formula otherwise. Numeral substitution gives a total recursive map
\[
 \sigma:\mathbb N^2\longrightarrow\mathbb N,
 \qquad \sigma(u,n)=\ulcorner\varphi_u(\overline n)\urcorner,
\]
where $\overline n$ is the numeral for $n$. To compute this map, first compute
the code of the closed term $\overline n$ by iteration on $n$. Traverse the
finite syntax tree of $\varphi_u$, replacing free occurrences of $v_0$ by
that term and retaining bound occurrences. The closed replacement term
introduces no variable capture. The constructors, parsing operations, and
traversal are effective, so this is a terminating algorithm jointly in $u$
and $n$. Its output is a sentence code, giving $s\sigma=\sigma$.

An arithmetical theory $T$ \emph{strongly defines} $A\subseteq\mathbb N$
\emph{numeralwise} by $\varphi(v_0)$ if
\[
 n\in A\ \Longrightarrow\ T\vdash\varphi(\overline n),
 \qquad
 n\notin A\ \Longrightarrow\ T\vdash\neg\varphi(\overline n).
\]
If $T$ is consistent, neither opposite assertion is provable in the
respective case.

A complementary pair of c.e. domains on $\mathbb N$ consists of a recursive
set $A$ and its complement: dovetailing their enumerations decides
membership. Suppose $T$ is consistent and c.e., with theorem domain $t$,
and every recursive set is strongly definable numeralwise in $T$. For each
such $A$, choose an index $u$ of a representing formula and take the total
constant map $g(x)=u$. The two defining implications and consistency give
\[
 [\sigma(g\boxtimes1)]^*t=\pi_2^*1_A,
 \qquad
 [\sigma(g\boxtimes1)]^*t^-=\pi_2^*1_{\mathbb N\setminus A}.
\]
Thus $t$ satisfies the categorical representation condition.

\begin{corollary}[Arithmetic essential incompleteness]\label{cor:arithmetic}
Let $B$ be a consistent arithmetical theory in the effective language above.
Suppose every recursive subset of $\mathbb N$ is strongly definable
numeralwise in $B$. Then every consistent computably enumerable extension
$T$ of $B$ is incomplete. Here extension means that every theorem of $B$
is a theorem of $T$; no enumerability hypothesis is imposed on $B$.
\end{corollary}
\begin{proof}
Each formula strongly defining a recursive set in $B$ also strongly defines
it in $T$, because its positive and negative instances remain provable
under extension.
Consistency of $T$ gives the corresponding exact representation equations.
By \eqref{eq:arithmetic-closure}, its theorem domain $t$ belongs to $\Th_c$.
Theorem~\ref{thm:essential}, applied with base and extension both equal to
$t$, proves that $T$ is incomplete.
\end{proof}

\begin{thebibliography}{9}
\bibitem{DPH} R.~A.~Di Paola and A.~Heller,
Dominical categories: recursion theory without elements,
\emph{Journal of Symbolic Logic} \textbf{52}(3) (1987), 594--635.
\bibitem{Heller} A.~Heller,
An existence theorem for recursion categories,
\emph{Journal of Symbolic Logic} \textbf{55}(3) (1990), 1252--1268.
\bibitem{RestrictionI} J.~R.~B.~Cockett and S.~Lack,
Restriction categories I: categories of partial maps,
\emph{Theoretical Computer Science} \textbf{270}(1--2) (2002), 223--259.
\end{thebibliography}
\end{document}
